\documentclass[conference]{IEEEtran}
\usepackage{cite}
\usepackage{amsmath,amssymb,amsfonts}
\usepackage{algorithmic}
\usepackage{graphicx}
\usepackage{textcomp}
\usepackage{xcolor}
\usepackage{booktabs}
\usepackage{tabularx}
\usepackage{hyperref}
\usepackage{listings}

\hypersetup{
    colorlinks=true,
    linkcolor=blue,
    filecolor=magenta,      
    urlcolor=cyan,
    citecolor=blue
}

\begin{document}

\title{DrugInsight-RAG: A Comparative Study of Graph-Based and Vector-Based Retrieval-Augmented Generation for Drug Review Analysis}

\author{
\IEEEauthorblockN{Asfar Hossain Sitab}
\IEEEauthorblockA{\textit{Department of Computer Science and Engineering} \\
\textit{East West University}\\
Dhaka, Bangladesh \\
Student ID: 2026-2-74-008 \\
AIML506 --- Fall 2026}
\and
\IEEEauthorblockN{Parmita Hossain Simia}
\IEEEauthorblockA{\textit{Department of Computer Science and Engineering} \\
\textit{East West University}\\
Dhaka, Bangladesh \\
Student ID: 2026-2-74-009 \\
AIML506 --- Fall 2026}
}

\maketitle

\begingroup\renewcommand\thefootnote{}\footnote{Course Instructor: Dr. Mohammad Rifat Ahmmad Rashid. Assignment Deadline: 4 December 2026. Kaggle Notebook: \url{https://www.kaggle.com/code/asfarhossainsitab/aiml-506-project-drugreview}}\addtocounter{footnote}{-1}\endgroup

\begin{abstract}
This study compares a knowledge-graph retrieval-augmented generation (GraphRAG) system with a dense-vector retrieval-augmented generation (Vector RAG) system for answering research questions about patient-authored drug reviews. The executed Kaggle notebook records 392,510 source records from the DrugRev dataset and a fixed-seed, ten-condition sample of 4,500 reviews. Both systems operate on 5,584 common evidence records: Neo4j Aura uses a six-label, seven-relation property graph (8,411 nodes, 20,599 relationships) and parameterized Cypher retrieval, whereas ChromaDB uses 384-dimensional MiniLM embeddings and cosine retrieval. Mistral 7B Instruct v0.3 is held constant as the quantized primary generator. Across 30 notebook-generated questions (six each of five types), graph retrieval achieved Hit@5 = 1.000 against 0.333 for vector retrieval, MRR = 0.967 against 0.130, and key-item accuracy = 0.667 against 0.633. Vector retrieval was faster, with median end-to-end latency of 7.07 s versus 8.95 s. A four-model ablation and five case analyses identify retrieval, generation, and benchmark-design failures, including a documented ten-way tie on aggregation queries.
\end{abstract}

\begin{IEEEkeywords}
GraphRAG, vector RAG, retrieval-augmented generation, Neo4j, ChromaDB, drug reviews, evidence grounding, evaluation, large language models.
\end{IEEEkeywords}

\section{Introduction}
Patient-authored medication reviews are an accessible source of experiential narratives and associations among drug names, indications, subjective ratings, duration of intake, and adverse-reaction words. An information system over these reviews should reliably answer bounded research questions: Which medications are reviewed for a condition? Which reviews mention a particular symptom? How frequently does a phrase appear? These questions require different forms of evidence and make the domain an ideal testbed for retrieval-augmented generation (RAG) \cite{lewis2020rag}.

RAG combines retrieved external evidence with a language model's generation process rather than relying solely on parametric memory \cite{guu2020realm, karpukhin2020dpr}. However, semantic proximity is not the same as relational or numerical correctness. A query about the most frequently reviewed medication or the intersection of conditions demands explicit grouping, counting, or joins; conversely, a query requesting the qualitative narrative of a review benefits from similarity-based text retrieval. Graph reasoning exposes typed entities and paths \cite{yasunaga2021qagnn}, while recent GraphRAG approaches emphasize structural organization for retrieval and synthesis \cite{edge2024graphrag}.

This study investigates whether structured, relationship-aware retrieval improves evidence selection, grounded answer quality, and reliability relative to dense-vector search on the same medication-review records.

\section{Dataset and Domain Understanding}
\subsection{Source and Sampling}
The dataset is \textbf{DrugRev: A comprehensive customers reviews on drugs purchasing and satisfaction}, Mendeley Data Version 1, DOI: 10.17632/64cc5w5dxy.1, under CC BY 4.0 \cite{rizwan2023drugrev}. The source shape is $392,510 \times 9$. The columns are \texttt{MedicineName}, \texttt{MedicineFor}, \texttt{ReviewDate}, \texttt{UserName}, \texttt{IntakeTime}, \texttt{Reviews}, \texttt{ReviewLength}, \texttt{Rating}, and \texttt{NumberOfLikes}.

A balanced sample of 4,500 reviews was drawn across 10 preselected clinical conditions: Depression, Anxiety, Pain, Insomnia, Type 2 Diabetes, High Blood Pressure, Bipolar Disorder, Weight Loss, Acne, and ADHD (seed 42, $n=450$ per condition). The sampled corpus covers 470 medications, 3,395 pseudonymous users, and 9 intake duration categories.

\subsection{Adverse-Reaction Extraction}
A 27-term dictionary covering clinical terms (nausea, headache, dizziness, fatigue, insomnia, weight gain, etc.) was matched using case-insensitive word-boundary regular expressions. The highest-frequency matched terms were insomnia (223 distinct reviews), nausea (187), weight loss (106), weight gain (87), and diarrhea (80).

\section{Methodology and Architecture}
\subsection{GraphRAG Pipeline (Neo4j Aura)}
The knowledge graph uses six project-prefixed labels to avoid collision: \texttt{AIML506Drug}, \texttt{AIML506Condition}, \texttt{AIML506Review}, \texttt{AIML506SideEffect}, \texttt{AIML506Patient}, and \texttt{AIML506IntakeDuration}. It incorporates \textbf{8,411 nodes} and \textbf{20,599 relationships} on Neo4j Aura (Instance02). Seven relationship types model domain semantics:
\begin{itemize}
    \item \texttt{POSTED} (4,500 edges): Patient to Review.
    \item \texttt{REVIEWS\_DRUG} (4,500 edges): Review to Drug.
    \item \texttt{FOR\_CONDITION} (4,500 edges): Review to Condition.
    \item \texttt{HAS\_INTAKE} (4,500 edges): Review to IntakeDuration.
    \item \texttt{REPORTS\_SIDE\_EFFECT} (1,317 edges): Review to SideEffect.
    \item \texttt{TREATS} (568 edges): Drug to Condition.
    \item \texttt{CAUSES\_SIDE\_EFFECT} (714 edges): Aggregated Drug to SideEffect co-mention association.
\end{itemize}
Entity linking is executed via boundary-matching dictionaries. A deterministic router selects one of five Cypher query paths: single-hop, multi-hop, aggregation, semantic-review lookup, and out-of-scope refusal.

\subsection{Vector RAG Pipeline (ChromaDB)}
All evidence records are serialized into plain-text documents with stable evidence IDs, giving Vector RAG access to the identical underlying information. ChromaDB is configured with cosine distance and \texttt{all-MiniLM-L6-v2} (384-dimensional dense embeddings). The collection indexes \textbf{5,584 records}: 4,500 review narratives, 568 medication-condition summaries, 470 medication summaries, 27 matched-term summaries, 10 condition summaries, and 9 intake duration summaries. Top $k=5$ documents are retrieved per query.

\subsection{Generator and Fairness Configuration}
The primary generator is \texttt{Mistral-7B-Instruct-v0.3} in 4-bit NF4 quantization on an NVIDIA Tesla T4 GPU. The system prompt instructs the model to answer strictly from provided evidence, avoid medical advice, treat symptoms as noncausal co-mentions, refuse unsupported questions, and cite exact bracketed evidence IDs. Decoding settings are identical for both systems: temperature 0, greedy decoding, max new tokens 160, input cap 3,072 tokens, and evidence budget 1,800 tokens.

\section{Empirical Evaluation and Results}
\subsection{Overall Benchmark Results (Table V)}
A 30-question evaluation benchmark was designed programmatically (6 single-hop, 6 multi-hop, 6 aggregation, 6 semantic/descriptive, and 6 unanswerable). Table \ref{tab:results} records the observed metrics from the executed Kaggle run.

\begin{table}[htbp]
\caption{Overall Observed Evaluation Metrics (Kaggle Executed Run)}
\label{tab:results}
\centering
\begin{tabularx}{\linewidth}{lrr}
\toprule
\textbf{Evaluation Metric} & \textbf{GraphRAG} & \textbf{Vector RAG} \\
\midrule
Hit@5 & \textbf{1.000} & 0.333 \\
Recall@5 & \textbf{1.000} & 0.333 \\
Precision@5 & \textbf{0.333} & 0.067 \\
Mean Reciprocal Rank (MRR) & \textbf{0.967} & 0.130 \\
nDCG@5 & \textbf{0.973} & 0.181 \\
\midrule
Key-item Accuracy & \textbf{0.667} & 0.633 \\
Token F1 (Overlap) & \textbf{0.410} & 0.366 \\
LLM-Judge Correctness (Qwen 2.5 7B) & 0.642 & \textbf{0.667} \\
Groundedness (LLM Judge) & 0.708 & \textbf{0.717} \\
Citation Validity & \textbf{0.833} & 0.250 \\
Out-of-Scope Refusal (6 items) & \textbf{1.000} & \textbf{1.000} \\
Graph Entity-Link F1 & \textbf{1.000} & N/A \\
\midrule
Median Total Time (s) & 8.948 & \textbf{7.069} \\
Median Retrieval Time (s) & 0.464 & \textbf{0.060} \\
Median Generation Time (s) & 8.392 & \textbf{6.996} \\
Mean Evidence Tokens & \textbf{385.8} & 463.3 \\
\bottomrule
\end{tabularx}
\end{table}

\subsection{Performance by Question Type (Table VI)}
Table \ref{tab:bytype} breaks down retrieval Hit@5 by query archetype ($n=6$ each).

\begin{table}[htbp]
\caption{Retrieval Hit@5 by Question Category}
\label{tab:bytype}
\centering
\begin{tabularx}{\linewidth}{lrr}
\toprule
\textbf{Category} & \textbf{GraphRAG} & \textbf{Vector RAG} \\
\midrule
Single-hop & \textbf{1.000} & 0.833 \\
Multi-hop & \textbf{1.000} & 0.000 \\
Aggregation & \textbf{1.000} & 0.333 \\
Semantic / Descriptive & \textbf{1.000} & 0.167 \\
Unanswerable & N/A & N/A \\
\bottomrule
\end{tabularx}
\end{table}

GraphRAG retrieved gold evidence in all answerable queries, while Vector RAG failed entirely on multi-hop questions (0.000) and struggled with aggregation (0.333).

\subsection{Four-Model Ablation Study (Table VII)}
Five cross-category queries were evaluated across four model architectures: Mistral 7B Instruct v0.3, Qwen 2.5 7B Instruct, Zephyr 7B Beta, and Phi-3.5 Mini Instruct (3.8B).

\begin{table}[htbp]
\caption{Four-Generator Model Ablation (5 Questions per Model)}
\label{tab:ablation}
\centering
\begin{tabularx}{\linewidth}{llrrr}
\toprule
\textbf{Model} & \textbf{System} & \textbf{Acc.} & \textbf{Token F1} & \textbf{Gen. Time (s)} \\
\midrule
Mistral 7B & Graph & 0.60 & 0.437 & 8.70 \\
Mistral 7B & Vector & 0.60 & 0.345 & 7.19 \\
Qwen 2.5 7B & Graph & 0.60 & 0.396 & 8.88 \\
Qwen 2.5 7B & Vector & 0.60 & 0.329 & 5.43 \\
Zephyr 7B & Graph & 0.60 & 0.447 & 12.19 \\
Zephyr 7B & Vector & 0.60 & 0.328 & 12.05 \\
Phi-3.5 Mini & Graph & 0.60 & 0.424 & 14.23 \\
Phi-3.5 Mini & Vector & 0.60 & 0.358 & 9.78 \\
\bottomrule
\end{tabularx}
\end{table}

\section{Discussion and Failure Analysis}
The notebook automatically identifies the five weakest observed question-system pairs:
\begin{enumerate}
    \item \textbf{AG1 / VectorRAG (Rec@5 = 0.0):} Failed set-level ranking of High Blood Pressure drugs; returned Toprol-XL and Bystolic instead of Azor, Atenolol, and Diovan due to vector chunking blindness.
    \item \textbf{AG4 / VectorRAG (Rec@5 = 0.0):} Global maximum miss; returned Orlistat (7 reviews) instead of Phentermine (100 reviews).
    \item \textbf{AG2 / VectorRAG (Rec@5 = 0.0):} Lexicon frequency miss; returned agitation (16 reviews) instead of insomnia (223), nausea (187), and weight loss (106).
    \item \textbf{MH6 / GraphRAG (Rec@5 = 1.0):} Generation/key-item mismatch; gold evidence was retrieved, but the generator failed to reproduce all required medication counts within the 160-token limit.
    \item \textbf{AG3 / VectorRAG (Rec@5 = 0.0):} Benchmark tie defect; Pain (450 reviews) tied with ADHD (450 reviews) under stratified sampling, but the gold label recognized only ADHD.
\end{enumerate}

\section{Conclusion}
GraphRAG achieved decisive advantages on evidence retrieval (Hit@5 of 1.000 vs 0.333, MRR of 0.967 vs 0.130) and citation validity (0.833 vs 0.250). Vector RAG provided faster execution (median total time 7.07 s vs 8.95 s) and slightly higher automated judge scores. For production clinical pharmacovigilance, a \textbf{Hybrid GraphRAG-VectorRAG} system represents the optimal architecture.

\begin{thebibliography}{00}
\bibitem{lewis2020rag} P. Lewis et al., ``Retrieval-augmented generation for knowledge-intensive NLP tasks,'' in \textit{Proc. NeurIPS}, 2020.
\bibitem{guu2020realm} K. Guu et al., ``REALM: Retrieval-augmented language model pre-training,'' in \textit{Proc. ICML}, 2020.
\bibitem{karpukhin2020dpr} V. Karpukhin et al., ``Dense passage retrieval for open-domain question answering,'' in \textit{Proc. EMNLP}, 2020.
\bibitem{yasunaga2021qagnn} M. Yasunaga et al., ``QA-GNN: Reasoning with language models and knowledge graphs for question answering,'' in \textit{Proc. NAACL}, 2021.
\bibitem{edge2024graphrag} D. Edge et al., ``From local to global: A graph RAG approach to query-focused summarization,'' \textit{arXiv:2404.16130}, 2024.
\bibitem{rizwan2023drugrev} M. Rizwan et al., ``DrugRev: A comprehensive customers reviews on drugs purchasing and satisfaction,'' \textit{Mendeley Data}, DOI: 10.17632/64cc5w5dxy.1, 2023.
\end{thebibliography}

\end{document}
